% 02-adding-problem.tex ;  The adding problem and results
%
% Scope: definition, why it requires memory, experimental setup, results

\section{Bài toán cộng}
\label{sec:ch02-adding}

\index{adding problem|(}%

\subsection{Định nghĩa}

Cho một chuỗi $T$ bước thời gian. Tại mỗi bước $t$, mạng nhận hai số:
\begin{itemize}
  \item Một giá trị ngẫu nhiên $v_t \sim \mathrm{Uniform}(-1, 1)$.
  \item Một cờ đánh dấu $m_t \in \{0, 1\}$. Cờ bằng $1$ tại \textbf{đúng hai}
    vị trí ngẫu nhiên trong chuỗi; bằng $0$ ở mọi vị trí còn lại.
\end{itemize}

Nhiệm vụ: tại bước cuối cùng $T$, mạng phải xuất ra tổng của hai giá trị tại
hai vị trí được đánh dấu.

\begin{example}
Với $T=10$, hai vị trí được đánh dấu là $t=3$ và $t=7$:
\[
\begin{array}{c|cccccccccc}
t      & 1 & 2 & 3    & 4 & 5 & 6 & 7    & 8 & 9 & 10 \\
\hline
v_t    & 0.5 & -0.3 & \mathbf{0.8} & 0.1 & -0.4 & 0.6 & \mathbf{0.2} & -0.7 & 0.9 & -0.1 \\
m_t    & 0 & 0 & \mathbf{1} & 0 & 0 & 0 & \mathbf{1} & 0 & 0 & 0
\end{array}
\]
Target: $0.8 + 0.2 = 1.0$.
\end{example}

Bài toán này buộc RNN phải làm hai việc: (a) bỏ qua phần lớn đầu vào (những
vị trí có $m_t = 0$ không mang thông tin hữu ích), và (b) nhớ giá trị tại vị
trí đánh dấu thứ nhất cho đến khi gặp vị trí thứ hai; có thể cách nhau hàng
chục bước.

\subsection{Kết quả}

Tôi huấn luyện RNN với $d_h = 24$, learning rate $0.01$, 5000 mẫu, khởi tạo
trọng số $\mathcal{N}(0, 0.1)$. Hàm mất mát là MSE trên output cuối cùng.

\begin{measured}{Adding problem, $d_h = 24$, 5000 mẫu}
\small
\begin{tabular}{@{}lrrrr@{}}
\toprule
$T$ & Loss trước & Loss sau & $\|\partial L / \partial \mat{W}_{hh}\|$ tại $t=1$ & tại $t=T$ \\
\midrule
10 & 2.86 & 0.67 & 0.387 & 0.228 \\
20 & 1.34 & 0.65 & 0.206 & 0.206 \\
50 & 0.002 & 0.66 & 0.235 & 0.196 \\
\bottomrule
\end{tabular}
\normalsize
\end{measured}

Vài điều đáng chú ý. Thứ nhất, loss giảm từ 2.86 xuống 0.67 ở $T=10$: RNN
học được bài toán ở chuỗi ngắn. Thứ hai, tỉ lệ norm gradient giữa $t=1$ và
$t=T$ gần bằng 1 ở mọi $T$; gradient chưa suy giảm đáng kể ở $d_h = 24$.
Điều này khớp với lý thuyết: vanishing gradient phụ thuộc vào $\|\mat{W}_{hh}\|$,
và với khởi tạo $\mathcal{N}(0, 0.1)$, ma trận $\mat{W}_{hh}$ có trị riêng
lớn nhất đủ nhỏ để gradient không bị triệt tiêu ở $T=50$.

Nhưng điều đó không có nghĩa là RNN \emph{học tốt} ở $T=50$. Loss sau huấn
luyện ở $T=50$ vẫn là 0.66, trong khi ở $T=10$ nó cũng là 0.67. Ở chuỗi dài
gấp 5 lần, RNN không tệ hơn; nhưng cũng không tốt hơn. Nó đã bão hòa.

Bài toán cộng không phải là bài kiểm tra khắc nghiệt nhất cho vanishing
gradient. Để thấy hiệu ứng rõ ràng hơn, tôi chuyển sang bài toán sao chép.

\index{adding problem|)}%
